\section{资本轰炸后的人为乱象}

上一章讲生产力，本章讲行业乱象。ChatGPT 之后，具身智能从小众概念变成资本宠儿，融资、估值和 demo 迅速升温。王鹤的批评核心是：行业最乱的地方在于分不清谁在创造生产力，谁只是在讲故事。更严重的是，有些团队展示时不明确说明遥操作，甚至把遥操作伪装成自主能力，这会直接破坏公众和客户对行业的信任。

这类问题不只是道德问题，也会变成技术问题。错误展示会把行业 reward model 带偏：团队为了融资和传播优化 demo，而不是为了真实场景成功率、成本下降和出货回流优化系统。长期下去，资本会惩罚整个行业，真正做生产力的团队也会被噪音拖累。

\begin{figure}[H]
\centering
\includegraphics[width=0.92\textwidth]{capital-chaos-reward.png}
\caption{资本轰炸后的乱象：谁创造生产力，谁只是在讲故事。自制概念图，依据 02:13:51--02:25:25 对谈内容与视频描述整理。}
\end{figure}

\begin{knowledgebox}{读图：错误 reward 会制造错误行为}
图中心是 reward。若 reward 来自融资额、估值和流量，团队会自然优化故事和演示；若 reward 来自真实展示、客户复购和万台应用，团队才会优化生产力。读图时要把“真实展示”看作行业基础设施：没有可信展示，市场无法判断进步。
\end{knowledgebox}

\subsection{五年万台应用的验证钟}

王鹤给出一个强判断：五年内领域必须出现万台以上的应用，如果做不到，具身智能这个领域会被证伪，至少当下叙事会被证伪。这个说法不一定是精确预测，但它提供了行业验收尺度。万台不是为了追求数字好看，而是为了证明硬件、软件、交付、维护、数据回流和客户价值已经跨过实验室阶段。

\begin{figure}[H]
\centering
\includegraphics[width=0.96\textwidth]{robotics-validation-clock.png}
\caption{五年验证时钟：五年内需要万台以上应用，否则领域叙事会被证伪。自制概念图，依据 02:13:51--02:25:25 对谈内容与视频描述整理。}
\end{figure}

\begin{importantbox}{真实展示原则}
可以遥操作，但必须说明遥操作；可以展示半成品，但必须说明边界；可以融资，但不能用不可复现的 demo 替代真实进展。对具身智能来说，行业信用本身就是一项公共资产。
\end{importantbox}

\begin{knowledgebox}{验证矩阵：什么才算真实进展}
\begin{tabularx}{\textwidth}{p{0.20\textwidth}p{0.34\textwidth}X}
\toprule
验证项 & 要看什么 & 不能用什么替代 \\
\midrule
公开展示 & 现场、无隐藏遥操作、可复现任务 & 剪辑视频和单次成功片段。 \\
场景运营 & 每天处理多少真实订单/任务，持续多久 & 战略合作和意向协议。 \\
出货规模 & 是否达到千台、万台级真实部署 & 工程样机数量。 \\
生产力 & 单位时间产出、故障率、人力兜底比例 & “像人”的外观和动作。 \\
数据飞轮 & 部署失败是否回流训练和产品改进 & 静态采集中心和一次性演示。 \\
\bottomrule
\end{tabularx}
\end{knowledgebox}

这个矩阵也解释了为什么王鹤反复强调不要砸行业招牌。中国面临老龄化和劳动力缺口，机器人生产力不是单纯创业风口，而可能关系到制造业和服务业的长期供给。如果行业被虚假展示和不切实际承诺拖入“冰河时代”，真正需要机器人补充劳动力的社会问题也会失去一个可能解法。

\subsection{黄仁勋插曲的意义}

前面讲的是行业信用，本节用访谈最后的插曲把技术生态收回来。黄仁勋来华访问时，王鹤和他同桌且挨着坐，观察到黄仁勋能吃辣、吃了不少水煮肉片，也对变脸表演给出热情反馈。这个片段看似八卦，但在笔记里可以作为一个侧面：NVIDIA、物理 AI、具身智能和中国机器人创业者已经处在同一个产业场域。算力、仿真、机器人本体和落地场景不再是平行线。

王鹤对这次同桌的解释也和技术主线有关：NVIDIA 很重视合成数据，因为合成数据需要 GPU 渲染、仿真和训练；如果这条路线走通，GPU 公司就能支撑具身智能的半边天。银河通用也向 NVIDIA 展示了用合成数据训练并迁移到真实世界的工作。这说明“生产力即产品”和“合成数据是数据 recipe 的关键”不是两个割裂判断，而是同一个闭环的两端：数据生产支撑能力，能力进入场景产生生产力，场景再回流真实数据。

\subsection{本章小结}

资本会放大机会，也会放大乱象。具身智能如果要建立长期信用，必须把 reward 从融资和 demo 转向真实展示、生产力、出货规模和客户价值。

